% 图 73.1 —— 由 `checks/emit_hand.py` 自动拼出（probe 精测 + prims2 图元分解）
%%
% ⚠ 这是**稿子不是成品**：跑 `checks/checkhand.py 73.1` 看指标 + 看叠合图。
%%
%% ⚠ 待核：
%   · 本图 probe 报出阴影族 1 个 —— **本脚本不生成阴影**（要照原书手写）；prims2 的 strokes 里可能含阴影线，核对时留意别画重
%   · 可疑字形 1 项（空心圈 ⊙ 常被认成字母 o）—— 见文件末
%%
\documentclass[12pt]{standalone}
\usepackage{fontspec}
\setsansfont{Arial}
\newfontfamily\mfallback{DejaVu Sans}
\usepackage{tikz}
\begin{document}
\begin{tikzpicture}[x=1pt, y=-1pt, line width=7.4pt, line cap=round, line join=round]

  %% ════ 填充区（prims2 的 fills）════

  %% ════ 直线段（prims2 的 strokes）════
  \draw[line width=7.4pt] (61.0,127.0) -- (67.0,131.0);
  \draw[line width=5.0pt] (60.0,130.0) -- (60.0,291.0);
  \draw[line width=5.0pt] (875.1,138.9) -- (885.7,126.3);
  \draw[line width=5.0pt] (885.7,126.3) -- (884.0,109.0);
  \draw[line width=5.0pt] (545.0,146.5) -- (549.0,120.1);
  \draw[line width=5.0pt] (549.0,120.1) -- (563.1,93.9);
  \draw[line width=6.0pt] (745.8,269.0) -- (728.0,266.0);
  \draw[line width=4.0pt] (319.0,131.0) -- (317.7,290.3);
  \draw[line width=5.0pt] (317.7,290.3) -- (167.0,293.0);
  \draw[line width=4.0pt] (882.0,187.0) -- (886.0,187.0);
  \draw[line width=5.0pt] (882.0,187.0) -- (877.0,192.0);
  \draw[line width=6.1pt] (59.0,129.0) -- (54.0,132.0);
  \draw[line width=5.0pt] (165.0,292.0) -- (61.0,292.0);
  \draw[line width=6.5pt] (312.0,133.0) -- (318.0,130.0);
  \draw[line width=5.0pt] (652.0,103.0) -- (644.0,116.8);
  \draw[line width=5.5pt] (189.0,21.0) -- (185.0,27.0);
  \draw[line width=3.0pt] (734.0,356.0) -- (719.0,356.0);
  \draw[line width=5.0pt] (922.9,147.1) -- (922.0,119.2);
  \draw[line width=4.5pt] (922.0,119.2) -- (904.5,95.5);
  \draw[line width=5.0pt] (646.1,78.0) -- (613.3,102.7);
  \draw[line width=9.7pt] (886.0,187.0) -- (886.0,196.0);
  \draw[line width=7.7pt] (487.0,118.0) -- (489.0,125.0);
  \draw[line width=5.0pt] (710.0,212.0) -- (716.0,249.0);
  \draw[line width=4.0pt] (727.0,265.0) -- (719.0,251.0);
  \draw[line width=5.0pt] (659.0,129.0) -- (670.9,118.1);
  \draw[line width=5.0pt] (857.1,115.1) -- (872.0,125.0);
  \draw[line width=5.0pt] (186.0,19.0) -- (60.7,20.3);
  \draw[line width=4.0pt] (60.7,20.3) -- (61.0,126.0);
  \draw[line width=6.0pt] (505.0,130.0) -- (490.0,126.0);
  \draw[line width=5.0pt] (190.0,20.0) -- (317.3,19.3);
  \draw[line width=5.0pt] (317.3,19.3) -- (319.0,128.0);
  \draw[line width=3.0pt] (488.0,126.0) -- (485.8,124.0);
  \draw[line width=6.0pt] (712.0,255.0) -- (715.0,250.0);
  \draw[line width=6.5pt] (166.0,294.0) -- (162.0,299.0);

  %% ════ 曲线（prims2 的 curves —— 三段式，坐标同为 (y,x)）════
  \draw[line width=5.0pt] (732.0,162.0) .. controls (720.7,175.8) and (709.3,190.9) .. (711.0,210.0);
  \draw[line width=5.0pt] (733.0,161.0) .. controls (781.4,159.3) and (832.7,162.0) .. (875.1,138.9);
  \draw[line width=5.0pt] (727.0,266.0) .. controls (655.4,260.6) and (557.8,230.3) .. (545.0,146.5);
  \draw[line width=5.0pt] (563.1,93.9) .. controls (629.6,18.9) and (737.5,11.1) .. (832.1,25.1);
  \draw[line width=4.0pt] (832.1,25.1) .. controls (898.7,36.6) and (998.9,83.8) .. (977.0,168.0);
  \draw[line width=5.0pt] (977.0,168.0) .. controls (936.2,252.4) and (832.4,268.6) .. (745.8,269.0);
  \draw[line width=5.0pt] (877.0,192.0) .. controls (826.5,214.2) and (767.1,221.4) .. (711.0,211.0);
  \draw[line width=5.0pt] (644.0,116.8) .. controls (653.9,153.9) and (701.5,154.0) .. (732.0,161.0);
  \draw[line width=5.0pt] (886.0,187.0) .. controls (901.1,176.9) and (916.7,164.6) .. (922.9,147.1);
  \draw[line width=5.0pt] (904.5,95.5) .. controls (832.1,43.2) and (725.9,42.2) .. (646.1,78.0);
  \draw[line width=5.0pt] (613.3,102.7) .. controls (602.7,115.2) and (599.5,132.3) .. (603.0,148.6);
  \draw[line width=4.0pt] (603.0,148.6) .. controls (622.5,190.1) and (670.1,202.5) .. (710.0,211.0);
  \draw[line width=5.0pt] (670.9,118.1) .. controls (727.2,90.2) and (799.3,93.4) .. (857.1,115.1);
  \draw[line width=4.0pt] (485.8,124.0) .. controls (451.8,120.8) and (415.6,119.6) .. (382.0,125.0);

  %% ════ 实心点 ════
  \fill (324.0,128.5) circle (5.6);
  \fill (487.5,130.5) circle (4.5);
  \fill (721.0,247.5) circle (4.9);
  \fill (59.2,292.4) circle (6.8);

  %% ════ 标签（probe 直接给的 \node 行）════
  %% ⚠ 全部补了 fill=white：文字在最上层，否则与曲线交叉干扰
     \node[anchor=center,inner sep=0pt,fill=white,font=\fontsize{32.0}{40.0}\selectfont] at (192.1,53.1) {\textsf{\textbf{\textit{a}}}};   % box(183,43,16,17)
     \node[anchor=center,inner sep=0pt,fill=white,font=\fontsize{24.5}{30.6}\selectfont] at (205.9,62.9) {\textsf{\textbf{1}}};   % box(199,54,8,17)
     \node[anchor=center,inner sep=0pt,fill=white,font=\fontsize{29.1}{36.4}\selectfont] at (434.2,79.2) {\textsf{\textbf{\textit{b}}}};   % box(425,67,16,22)
     \node[anchor=center,inner sep=0pt,fill=white,font=\fontsize{33.5}{41.9}\selectfont] at (879.0,103.3) {\textsf{\textbf{.}}};   % box(874,101,5,4)
     \node[anchor=center,inner sep=0pt,fill=white,font=\fontsize{29.8}{37.2}\selectfont] at (342.2,152.7) {\textsf{\textbf{\textit{b}}}};   % box(333,140,16,23)
     \node[anchor=center,inner sep=0pt,fill=white,font=\fontsize{29.1}{36.4}\selectfont] at (27.2,155.2) {\textsf{\textbf{\textit{b}}}};   % box(18,143,16,22)
     \node[anchor=center,inner sep=0pt,fill=white,font=\fontsize{23.7}{29.7}\selectfont] at (355.8,167.4) {\textsf{\textbf{1}}};   % box(349,159,8,16)
     \node[anchor=center,inner sep=0pt,fill=white,font=\fontsize{23.1}{28.9}\selectfont] at (41.4,169.2) {\textsf{\textbf{0}}};   % box(34,161,11,16)
     \node[anchor=center,inner sep=0pt,fill=white,font=\fontsize{32.0}{40.0}\selectfont] at (676.1,219.1) {\textsf{\textbf{\textit{a}}}};   % box(667,209,16,17)
     \node[anchor=center,inner sep=0pt,fill=white,font=\fontsize{32.0}{40.0}\selectfont] at (179.1,258.1) {\textsf{\textbf{\textit{a}}}};   % box(170,248,16,17)
     \node[anchor=center,inner sep=0pt,fill=white,font=\fontsize{24.1}{30.2}\selectfont] at (193.0,268.2) {\textsf{\textbf{0}}};   % box(185,260,12,16)
     \node[anchor=center,inner sep=0pt,fill=white,font=\fontsize{30.6}{38.3}\selectfont] at (37.5,317.5) {\textsf{\textbf{\textit{F}}}};   % box(28,305,19,23)
     \node[anchor=center,inner sep=0pt,fill=white,font=\fontsize{23.6}{29.6}\selectfont] at (197.1,342.9) {\textsf{\textbf{2}}};   % box(191,334,11,17)
     \node[anchor=center,inner sep=0pt,fill=white,font=\fontsize{23.7}{29.7}\selectfont] at (838.8,343.4) {\textsf{\textbf{1}}};   % box(832,335,8,16)
     \node[anchor=center,inner sep=0pt,fill=white,font=\fontsize{24.5}{30.6}\selectfont] at (772.9,344.9) {\textsf{\textbf{1}}};   % box(766,336,8,17)
     \node[anchor=center,inner sep=0pt,fill=white,font=\fontsize{23.7}{29.6}\selectfont] at (182.4,350.9) {\textsf{\textbf{\textit{f}}}};   % box(178,340,8,21)
     \node[anchor=center,inner sep=0pt,fill=white,font=\fontsize{29.6}{37.0}\selectfont] at (757.2,351.5) {\textsf{\textbf{\textit{B}}}};   % box(746,340,19,22)
     \node[anchor=center,inner sep=0pt,fill=white,font=\fontsize{31.5}{39.4}\selectfont] at (823.2,353.3) {\textsf{\textbf{\textit{S}}}};   % box(814,341,19,23)
     \node[anchor=center,inner sep=0pt,fill=white,font=\fontsize{36.3}{45.4}\selectfont] at (705.2,354.9) {\textsf{\textbf{\textit{x}}}};   % box(694,342,20,21)
     \node[anchor=center,inner sep=0pt,fill=white,font=\fontsize{37.1}{46.4}\selectfont] at (794.1,356.3) {\textsf{\textbf{×}}};   % box(785,345,17,17)
     \node[anchor=center,inner sep=0pt,fill=white,font=\fontsize{33.9}{42.4}\selectfont] at (735.7,355.6) {{\mfallback\textbf{−}}};   % box(719,348,17,5)

  % 钉死画布：否则超出的图元/控制点会把页面撑大，1:1 回比就失准
  \pgfresetboundingbox
  \path[use as bounding box] (0,0) rectangle (994,377);
\end{tikzpicture}
\end{document}

% ⚠ 待核清单（probe 拿不准，人眼过一遍）：
%   · 未识别/跳过：⑤ 跳过/未识别的字形
